\section{From Potential Outcomes to Other Causal Frameworks}\label{sec:extension}
We proposed a method based on weighted conformal inference which
produces interval estimates of counterfactuals and individual
treatment effects under the potential outcome framework. For
randomized experiments with perfect compliance, our method has
guaranteed coverage in finite samples without any modeling assumptions
on the data generating process. For randomized experiments with
ignorable compliance or general observational studies under the strong
ignorability assumption, our method is doubly robust in the sense that
the coverage is asymptotically guaranteed if either the conditional
quantiles of potential outcomes or the propensity scores are
consistently estimated. In contrast, existing methods may suffer from
a significant coverage deficit even in simple models. Furthermore, our
framework naturally extends to other populations by modifying the
weight function according to Table \ref{tab:weight_func}.
% \ejc{This
%   feels like a repetition but I think this is okay.}

The key observation is the covariate shift together with the
invariance of the conditional distribution: the observed distribution
of $(X, Y(1))$ is $P_{X\mid T = 1}\times P_{Y(1)\mid X}$ under
ignorability while the target distribution is
$P_{X}\times P_{Y(1)\mid X}$. Now the invariance of the conditional distribution is
also the enabling property in other frameworks of causal inference. As
a consequence, our method can be naturally extended to those settings,
and we close this paper by discussing two of them.

\subsection{Causal diagram framework}
Judea Pearl in his pioneering work \citep{pearl1995causal} introduced
a general framework of causal inference based on graphical models,
which has been widely applied, see
  e.g.~\cite{greenland1999causal, spirtes2000causation,
    richardson2013single, glymour2014discovering,
    tennant2019use}. This framework does not rely on the notion of
counterfactuals but instead defines causal effects through the
\texttt{do} operator, which modifies the observed distribution by
removing the causal paths that directly point to the intervention
variable. We refer the readers to \cite{pearl2018book} for the
philosophy and basics of causal diagrams and to \cite{pearl2016causal}
for the mathematical foundation of this framework.

We here only discuss the case where $T$ is the intervention variable,
$Y$ is the outcome variable and $X$ is a set of variables that
satisfies the back-door criterion. Roughly speaking, this means that $X$
includes all confounders and excludes all post-treatment variables. In
this case, the foundational result in \cite{pearl1995causal} shows
that
\[P_{(X, Y) \mid \mathrm{do}(T = t)} = P_{X}\times P_{Y\mid X, T = t}.\]
This is the target distribution to be inferred. In contrast, the observed distribution of $(X, Y)$ given $T = t$ is
\[P_{(X, Y)\mid T = t} = P_{X\mid T = t}\times P_{Y\mid X, T = t}.\]
Clearly, this has exactly the same structure as in the potential
outcome framework. As a consequence, weighted split-CQR can be applied
without any modification to produce doubly robust interval estimates
of $Y$ under the \texttt{do} intervention.


\subsection{Invariant prediction framework}

Invariant prediction is another framework proposed by
\cite{peters2016causal}. It is particularly powerful when there are
multiple data sources under different interventions, such as in gene
knockout experiments. Consider an outcome variable $Y$, a set of
interventions or covariates $X$ and an environment variable $E$ that
indicates the source of data. Then one basic setting under this
framework assumes that $Y\indep E \mid X$ while $X$ may depend on $E$ with
$X \mid E\sim P_{X}^{E}$. The goal is to predict the outcome under a
new environment. For simplicity, we assume that a dataset
$(X_{ij}, Y_{ij})_{i=1}^{n_{j}}$ is available for each environment
$e_{1}, \ldots, e_{J}$ and a testing dataset $(X_{i0})_{i=1}^{n_{0}}$
for the target environment $e_{0}$. Due to the invariance assumption,
on the $j$-th dataset, the observed distribution of $(X, Y)$ is
$P_{X}^{e_{j}}\times P_{Y\mid X}$ while the target distribution to be
inferred on is $P_{X}^{e_{0}}\times P_{Y\mid X}$.

Again, this is reduced to a problem of constructing valid prediction
intervals under covariate shifts. When $J = 1$, this has exactly the
same structure as in the potential outcome framework and thus weighted
split-CQR with weight function $dP_{X}^{e_{0}}(x) / dP_{X}^{e_{1}}(x)$
produces doubly robust intervals of $Y$ under environment $e_{0}$. When $J > 1$, we can in principle apply the general
weighted conformal inference techniques from
\cite{barber2019conformal}. However the weight function becomes much
more complicated than that in Algorithm \ref{algo:split-CQR}. An
alternative approach is to create a weighted population from
environments $e_{1}, \ldots, e_{J}$ with observed distribution
\[\lb\sum_{j=1}^{J}q_{j}P_{X}^{e_{j}}\rb\times P_{Y\mid X},\]
and apply Algorithm \ref{algo:split-CQR} on this pseudo dataset with weight $w(x) = 1 / \sum_{j=1}^{J}q_{j}(dP_{X}^{e_{j}} / dP_{X}^{e_{0}})(x)$. The weights $q_{j}$ can be chosen through certain balancing procedures that forces the covariate distribution to approximate $P_{X}^{e_{0}}$. We leave the formal development of this idea to future work.
%%% Local Variables:
%%% mode: latex
%%% TeX-master: t
%%% End:
